
\documentclass[10pt, a4paper]{article}
\usepackage[utf8]{inputenc}
\usepackage{amsmath}
\usepackage{geometry}
\geometry{a4paper, margin=0.8in}
\usepackage{amssymb}
\usepackage{breqn}
\usepackage{booktabs}
\usepackage{url}

% Ajustes para permitir ecuaciones gigantes
\allowdisplaybreaks
\setlength{\jot}{8pt}

\title{Análisis Matemático de Algoritmo \\ \large Generado por Project Diophantus}
\author{}
\date{\today}

\begin{document}
\maketitle
\section*{Resumen Ejecutivo}
Este documento presenta la traducción completa de un algoritmo computacional a tres formas matemáticas distintas.
Las ecuaciones resultantes, debido a la complejidad logarítmica del algoritmo original (Miller-Rabin), son extensas pero polinomialmente acotadas.
\part{Función de Transición}
\begin{align*}
n[t+1] &= n \\
result[t+1] &= 0 \\
\end{align*}\part{Conversión a Polinomio}
\subsection*{Ecuación Unificada P=0}
\begin{align*}
& (n[t+1] - (n))^{2} + (result[t+1] - (0))^{2} = 0
\end{align*}\part{Fórmulas Generadoras}
\begin{small}
\begin{align*}
G &= 1 - ((n[t+1] - (n))^{2} + (result[t+1] - (0))^{2}) \\
\end{align*}
\end{small}\part{Fórmulas Matemáticas (Estilo Willans)}
\section{Representación Simbólica}
\begin{dmath*}
f(N) &= \sum_{R=0}^{\infty} R \cdot \lfloor \frac{1}{1 + \mathcal{D}_{sym}} \rfloor \\
\mathcal{D}_{sym} &= (\mathit{n} - (\mathit{n}))^2 + (R - (0))^2
\end{dmath*}\newpage
\section{Representación Operacional}
\begin{dmath*}
f(N) &= \sum_{R=0}^{\infty} R \cdot \left\lfloor \frac{1}{1 + \mathcal{D}(N, R, \mathbf{x})} \right\rfloor \\
\mathcal{D}(N, R, \mathbf{x}) &= \sum_{\mathbf{x} \in \mathbb{N}^{1}} \left( (x_{1} - (x_{1}))^2 + (R - (0))^2 \right)
\end{dmath*}
Vector de estado: $\mathbf{x} \in \mathbb{R}^{1}$ (variables internas anonimizadas).\\\part{Sistema Dinámico (Recurrencia Infinita)}
\begin{align*}
Transición de estado x(t+1) &= F(x(t)). Las C_n son subexpresiones \\
\text{comunes reutilizadas; las P_f, relaciones de funciones recursivas.} \\
\text{# Recurrencias de estado} \\
n(t+1) &= n \\
result(t+1) &= 0 \\
\end{align*}\end{document}